\documentclass[11pt,a4paper]{article}

\usepackage[utf8]{inputenc}
\usepackage[T1]{fontenc}
\usepackage{amsmath,amssymb,amsthm}
\usepackage{graphicx}
\usepackage{float}
\usepackage{booktabs}
\usepackage{geometry}
\usepackage{hyperref}

\geometry{margin=2.5cm}

% Theorem environments
\theoremstyle{plain}
\newtheorem{theorem}{Theorem}[section]
\newtheorem{proposition}[theorem]{Proposition}
\newtheorem{lemma}[theorem]{Lemma}
\newtheorem{corollary}[theorem]{Corollary}

\theoremstyle{definition}
\newtheorem{definition}[theorem]{Definition}
\newtheorem{example}[theorem]{Example}

\theoremstyle{remark}
\newtheorem{remark}[theorem]{Remark}

\title{Mathematical Report Title}
\author{Author Name}
\date{\today}

\begin{document}

\maketitle

\begin{abstract}
% Concisely summarize the problem, method, and main results.
\end{abstract}

\section{Introduction}
% Provide motivation, define non-standard notation, and state the objective.

\section{Main Results}

\begin{definition}
% State a definition needed by the results.
\end{definition}

\begin{theorem}
% State the primary mathematical claim clearly, including all hypotheses.
\end{theorem}

\begin{proof}
% Give a logically complete proof. Explain non-obvious transformations.
\end{proof}

\section{Computational Verification and Figures}

% Distinguish exact symbolic calculations, numerical approximations, and
% empirical observations. Do not call numerical agreement a proof.

% Example figure block (uncomment when an actual file is available):
% \begin{figure}[ht]
% \centering
% \includegraphics[width=0.75\textwidth]{figures/example.pdf}
% \caption{A meaningful description of the figure.}
% \label{fig:results}
% \end{figure}

\section{Conclusion}
% Summarize the proved results, computational evidence, limitations, and open questions.

\end{document}
